\documentclass[conference]{IEEEtran}
\IEEEoverridecommandlockouts
% The preceding line is only needed to identify funding in the first footnote. If that is unneeded, please comment it out.
\usepackage{cite}
\usepackage{amsmath,amssymb,amsfonts}
\usepackage{algorithmic}
\usepackage{graphicx}
\usepackage{textcomp}
\usepackage{xcolor}
\usepackage{minted}
\def\BibTeX{{\rm B\kern-.05em{\sc i\kern-.025em b}\kern-.08em
    T\kern-.1667em\lower.7ex\hbox{E}\kern-.125emX}}
\begin{document}

\title{Human Faces Generation with Diffusion Models\\
}

\author{
\IEEEauthorblockN{Zain Abideen\textsuperscript{1},
Hakan Demirer\textsuperscript{1},
Matthew Osborne\textsuperscript{1},
Kevin O'Shaughnessy\textsuperscript{1},
Pranav Pokhrel\textsuperscript{1}}
\IEEEauthorblockA{\textsuperscript{1}\textit{University of Surrey}\\
}
}


\maketitle

\begin{abstract}
Generative modelling aims to learn complex data distributions in order to generate realistic samples. In this project, we implement a diffusion-based model for unconditional human face generation using the CelebA-HQ dataset. The model is based on the Denoising Diffusion Probabilistic Model framework, where a neural network is trained to reverse a gradual Gaussian noising process. Experiments were conducted using a restricted dataset of 2700 training images and 300 test images, with performance evaluated using the Fréchet Inception Distance. We systematically investigated the effects of image resolution, model capacity, training duration, learning rate, and stabilisation techniques. Results show that longer training substantially improves generation quality, but also introduces instability when no stabilisation is used. Exponential Moving Average improves robustness and sample quality, reducing FID from approximately 234 to 173 in an early controlled comparison, and enabling continued improvement over longer training. The best results were obtained by combining EMA with cosine learning rate scheduling, achieving a final FID of 41.82 after 600 epochs. Qualitative results also showed clearer facial structure, improved symmetry, and reduced artifacts. Overall, the project demonstrates that diffusion models can generate realistic human faces even under limited data and computational constraints, with training stability playing a key role in final performance.
\end{abstract}

\begin{IEEEkeywords}
component, formatting, style, styling, insert
\end{IEEEkeywords}

\section{Introduction}
Generative modelling aims to learn data distributions in order to generate realistic samples. In image generation, this involves producing images that are visually indistinguishable from real data, with applications in content creation and data augmentation.
Human face generation is a challenging task due to complex structures and fine details, making it a strong benchmark for generative models. While Generative Adversarial Networks (GANs) \cite{b6} have achieved strong results, they often suffer from instability and mode collapse.
Diffusion models provide a more stable alternative by learning to iteratively transform noise into structured images through a denoising process. This approach has achieved state-of-the-art performance across multiple domains \cite{b9}.
In this work, we implement a diffusion-based model for human face generation using the CelebA-HQ dataset. We investigate the effects of architecture, resolution, training duration, and stabilization techniques such as Exponential Moving Average (EMA). Performance is evaluated using the Fréchet Inception Distance (FID), with significant improvements achieved through systematic experimentation.

\section{Literature Review}
Generative modeling aims to learn a probability distribution over high-dimensional data such that new samples can be drawn from it. Image synthesis has long served as a proving ground for generative methods, and for much of the last decade the dominant approach was Generative Adversarial Networks (GANs) \cite{b6}. GANs frame image generation as a two-player game between a generator that produces candidate samples from random noise and a discriminator that distinguishes them from real data.

Despite producing visually impressive results, GANs are notoriously difficult to train: the adversarial objective is prone to instability, mode collapse reduces the diversity of generated samples, and convergence is sensitive to architectural and hyperparameter choices. These limitations motivated the search for generative models with more stable training dynamics and stronger likelihood-based foundations.

Diffusion models provide such an alternative. The core idea, first proposed by Sohl-Dickstein et al. [4] and drawn from nonequilibrium thermodynamics, is to define a forward process that gradually corrupts data by adding Gaussian noise over many timesteps, and then to learn a reverse process that removes noise step by step, recovering samples from the data distribution.

Ho et al. [3] made this framework practical with Denoising Diffusion Probabilistic Models (DDPMs), and this formulation of DDPM forms the basis of the model implemented in this project. Song et al. [2] unified discrete-time denoising diffusion with continuous-time score-based generative modeling under a single stochastic differential equation framework, showing that both approaches can be viewed as learning the score function of a data distribution perturbed by noise at varying scales. 

Shortly afterwards, Dhariwal and Nichol [1] demonstrated empirically that diffusion models outperform GANs on standard image synthesis benchmarks across multiple resolutions, through a combination of architectural improvements, a learned reverse-process variance, and classifier-based guidance for conditional generation.

Subsequent work has identified several refinements to the original DDPM formulation that improve sample quality at modest implementation cost. Nichol and Dhariwal [8] introduced a cosine noise schedule in which the cumulative signal-preservation factor decays as a shifted squared cosine. The cosine schedule destroys image information more gradually across timesteps, avoiding the rapid saturation observed at late timesteps under linear scheduling, and yields improved FID at equivalent training cost. 

Dhariwal and Nichol [1] further showed that adaptive group normalization — in which the affine parameters of group normalization are predicted from the timestep embedding rather than learned as constants — provides a stronger form of timestep conditioning than the additive injection used in the original DDPM. The use of exponential moving averages of model parameters during sampling, while not introduced by any single paper, has become standard practice across the diffusion literature for stabilizing sample quality. 

Progress in generative image modeling depends on standardized benchmarks and quantitative evaluation. For unconditional face generation, the CelebA-HQ dataset introduced by Karras et al. [7] has become a widely adopted benchmark, providing 30,000 high-quality face images and supporting evaluation across GANs, score-based approaches, and diffusion models. 

To assess sample quality, the Fréchet Inception Distance (FID) proposed by Heusel et al. [5] has emerged as the standard metric. FID embeds both real and generated images into the feature space of a pretrained Inception-v3 network, fits a multivariate Gaussian to each set of features, and computes the Fréchet distance between the two Gaussians. 

\section{Methodology}

\subsection{Diffusion Model Framework}
This work is based on the Denoising Diffusion Probabilistic Model (DDPM) framework. Diffusion models consist of a forward process and a reverse process. The forward process gradually adds Gaussian noise to an image over a fixed number of timesteps until it becomes pure noise. The reverse process is learned by a neural network, which iteratively removes noise to reconstruct the original image.
The model is trained to predict the noise added at each timestep. Starting from random noise, the learned reverse process enables generation of realistic images through successive denoising steps.

\subsection{Dataset and Preprocessing}
Experiments were conducted on the CelebA-HQ dataset, restricted to 2700 training images and 300 test images as specified.
Images were resized to the required resolution (64×64 or 128×128), and normalized for stable training. No additional augmentation was applied to ensure consistency across experiments.

\subsection{Training Setup}
Models were trained using the Adam optimizer with a learning rate of 0.0002 and a batch size of 1. The diffusion process used 1000 timesteps.
Key variables explored include:

\begin{itemize}
  \item image resolution
  \item model capacity (base channels)
  \item number of training epochs
\end{itemize}

\subsection{Stabilization Techniques}
To improve training stability, Exponential Moving Average (EMA) was applied to model parameters, maintaining a smoothed version of the weights during training. EMA reduces noisy updates and improves sample quality during generation.
Additionally, gradient clipping was used to limit large gradient updates and further stabilize training.

\subsection{Evaluation Metric}
Performance was evaluated using the Fréchet Inception Distance (FID), which measures the similarity between real and generated image distributions. Lower FID indicates better sample quality and diversity.

\section{Experiments}
All experiments were conducted using a Denoising Diffusion Probabilistic Model (DDPM) trained on the CelebA-HQ dataset. In accordance with the coursework specification, the dataset was limited to 2700 training images and 300 test images. Model performance was evaluated using the Fréchet Inception Distance (FID), computed between 300 generated images and the test set.
The baseline configuration was defined as:
\begin{itemize}
  \item Image resolution: 128 × 128,
  \item Base channels: 64
  \item Batch size: 1
  \item Learning rate: 0.0002
  \item Diffusion timesteps: 1000
\end{itemize}

Experiments were designed to systematically explore the impact of:

\begin{enumerate}
  \item Model capacity
  \item Image resolution
  \item Training duration
  \item Learning rate
  \item Training stabilization techniques
\end{enumerate}

\subsection{Baseline Performance}
The baseline experiment (Experiment B) established a reference point using the default configuration. While it produced visually plausible face structures, the FID score remained relatively high, indicating limited sample quality and diversity.
This baseline provided a foundation for subsequent improvements through architectural and training modifications.

\subsection{Effect of Model Capacity}
To evaluate the effect of model capacity, experiments were conducted with varying base channel sizes.
Reducing the number of channels (Experiment C, ch32) resulted in poorer performance, indicating that the model lacked sufficient capacity to capture the complexity of facial features. Conversely, increasing capacity (Experiment G, ch96) did not yield improvements and instead led to worse FID scores, suggesting diminishing returns and potential optimization instability.
These results indicate that a moderate model size (ch64) provides the best balance between representational power and training stability.

\subsection{Effect of Image Resolution}
The effect of image resolution was investigated by comparing 128×128 and 64×64 configurations.
Lower resolution models (Experiment D) achieved relatively competitive FID scores; however, visual inspection revealed reduced detail and realism in generated faces. While lower resolution simplifies the learning problem, it limits the model’s ability to generate high-quality images.
Overall, 128×128 resolution provided better qualitative results and improved scalability, and was therefore retained for subsequent experiments.

\subsection{Effect of Training Duration}
Training duration was found to have a significant impact on model performance.
Increasing the number of epochs led to consistent improvements in FID:

\begin{itemize}
\item 100 epochs $\to$ FID $\approx$ 181
\item 150 epochs $\to$ FID $\approx$ 122
\item 200 epochs $\to$ FID $\approx$ 118
\end{itemize}

However, extending training further without additional stabilization led to degradation:

\begin{itemize}
\item 250 epochs $\to$ FID $\approx$ 159
\end{itemize}

This indicates that while longer training improves performance, overtraining can negatively affect sample quality, likely due to instability in later training stages.

\subsection{Effect of Learning Rate}
A lower learning rate (Experiment J, 1e-4) was also evaluated. While this configuration resulted in more gradual training, it led to slower convergence and inferior FID scores compared to the baseline learning rate of 0.0002.
Therefore, 0.0002 was retained as the optimal learning rate, providing a better balance between convergence speed and stability.

\subsection{EMA for Training Stabilization}
To address instability in longer training runs, Exponential Moving Average (EMA) was introduced. EMA maintains a smoothed version of model parameters, reducing the impact of noisy updates during optimization.
A controlled experiment was conducted to validate its effectiveness over 10 epochs:

\begin{itemize}
\item Raw model $\to$ FID $\approx$ 234
\item EMA model $\to$ FID $\approx$ 173
\end{itemize}

This demonstrated a substantial improvement even in early training.
EMA was then applied to longer training runs, producing significant gains:

\begin{itemize}
\item 200 epochs (EMA) $\to$ FID $\approx$ 95
\item 250 epochs (EMA) $\to$ FID $\approx$ 76
\item 300 epochs (EMA) $\to$ FID $\approx$ 62
\end{itemize}

Notably, EMA prevented the degradation observed in non-EMA runs and enabled continued improvements with extended training.
Qualitative Results
Generated samples show a clear progression in visual quality as training improves. Early-stage models produce noisy or distorted faces, while later models capture more realistic facial structures, including symmetry, texture, and finer details.
EMA-based models exhibit:
- smoother outputs
- fewer artifacts
- more consistent facial features

A comparison between generated and real images demonstrates that the final model closely approximates the distribution of real faces.

\subsection{Results Summary}
Table 1 summarizes the most important experimental results.

\begin{table}[htbp]
\caption{Summary of Experiments}
\begin{center}
\begin{tabular}{|c|c|c|}
\hline
\textbf{Experiment} & \textbf{\textit{Configuration}}& \textbf{\textit{FID}} \\
\hline
Baseline (B) &  128px, ch64, 20 epochs & 309.14 \\
C & 128px, ch32 & 261.26 \\
D & 64px, ch64 & 216.12 \\
E & 128px, ch64, 50 epochs & 282.86 \\
F100 & 128px, ch64, 100 epochs & 181.60 \\
G & 128px, ch96 & 247.07 \\
H & ch32, 50 epochs & 209.64 \\
I & 128px, ch96, batch 1, 50 epochs & 257.59 \\
K & 128px, ch96, lr=1e-4, 20 epochs & 247.59 \\
F150 & 128px, ch64, 150 epochs & 122.23 \\
F200 & 128px, ch64, 200 epochs & 118.30 \\
F250 & 128px, ch64, 250 epochs & 159.40 \\
EMA raw check & 128px, ch64, 10 epochs & 234 \\
EMA (10 epochs) & + EMA & 173 \\
F200 EMA & 128px, ch64, EMA & ~95 \\
F250 EMA & 128px, ch64, EMA & ~76 \\
F300 EMA & 128px, ch64, EMA & ~62 \\
F350 EMA + cos LR & 128px, ch64, EMA, cos LR, 350 epochs & 45.02 / 46.73 \\
F400 EMA + cos LR & 
128px, ch64, EMA, cos LR, 400 epochs &
44.65 \\
F450 EMA + cos LR &
128px, ch64, EMA, cos LR, 450 epochs &
42.64 \\
F500 EMA + cos LR &
128px, ch64, EMA, cos LR, 500 epochs &
45.92 \\
F550 EMA + cos LR &
128px, ch64, EMA, cos LR, 550 epochs &
43.90 \\
F600 EMA + cos LR &
128px, ch64, EMA, cos LR, 600 epochs &
41.82 \\
F300 EMA + cos noise &
128px, ch64, EMA, cos noise &
314 \\
Warmup EMA 100 &
ch128, bs16, warmup, EMA, 100 epochs &
108.3 \\
Warmup EMA 300 &
ch128, bs16, warmup, EMA, 300 epochs &
Running \\
Warmup EMA 400 &
ch128, bs16, warmup, EMA, 400 epochs &
Planned \\
J100 & 128px, ch64, lr=1e-4, 100 epochs & 240.12 \\
J250 &  & \\

M10 & EMA and AdaGN, bs=4, 10 epochs & 238.80 \\
EMA only &  10 epochs, bs=1 & 278.61 \\
EMA w accum steps & 100 epochs, bs=4*16 accum step, LR/4 & 236.64 \\
clamp to [-1, 1] & as above w prediction step + clamp to [-1, 1] & 149.87 \\
clamp to [-1, 1] ddim sample & as above w 50 ddim sampling, eta=0  & 149.82 \\
clamp to [-1, 1] ddim sample & as above w 100 ddim sampling, eta=0 & 162.76 \\
clamp to [-1, 1] ddim sample & as above w 250 ddim sampling, eta=0 & 160 \\
clamp to [-1, 1] ddim sample & as above w 1000 ddim sampling, eta=0 & 173.24 \\

EMA + zero init conv & 10 epochs, bs=1 & ? \\
EMA only & 20 epochs, bs=1 & 299.14 \\
M200 & EMA and AdaGN, bs=4, 200 epochs & 161.31 \\

\hline
\end{tabular}
\label{tab1}
\end{center}
\end{table}


\subsection{U-Net}\label{AA}
Define abbreviations and acronyms the first time they are used in the text, 
even after they have been defined in the abstract. Abbreviations such as 
IEEE, SI, MKS, CGS, ac, dc, and rms do not have to be defined. Do not use 
abbreviations in the title or heads unless they are unavoidable.

\subsection{Training setup}

\subsection{Hyperparameters}

\subsection{Data preprocessing}

\begin{itemize}
\item Use either SI (MKS) or CGS as primary units. (SI units are encouraged.) English units may be used as secondary units (in parentheses). An exception would be the use of English units as identifiers in trade, such as ``3.5-inch disk drive''.
\item Avoid combining SI and CGS units, such as current in amperes and magnetic field in oersteds. This often leads to confusion because equations do not balance dimensionally. If you must use mixed units, clearly state the units for each quantity that you use in an equation.
\item Do not mix complete spellings and abbreviations of units: ``Wb/m\textsuperscript{2}'' or ``webers per square meter'', not ``webers/m\textsuperscript{2}''. Spell out units when they appear in text: ``. . . a few henries'', not ``. . . a few H''.
\item Use a zero before decimal points: ``0.25'', not ``.25''. Use ``cm\textsuperscript{3}'', not ``cc''.)
\end{itemize}


\subsection{Generated image examples}

\subsection{Training loss curves}

\subsection{Equations}
Number equations consecutively. To make your 
equations more compact, you may use the solidus (~/~), the exp function, or 
appropriate exponents. Italicize Roman symbols for quantities and variables, 
but not Greek symbols. Use a long dash rather than a hyphen for a minus 
sign. Punctuate equations with commas or periods when they are part of a 
sentence, as in:
\begin{equation}
a+b=\gamma\label{eq}
\end{equation}

Be sure that the 
symbols in your equation have been defined before or immediately following 
the equation. Use ``\eqref{eq}'', not ``Eq.~\eqref{eq}'' or ``equation \eqref{eq}'', except at 
the beginning of a sentence: ``Equation \eqref{eq} is . . .''

Only include the most important 2-3 figures in the main body.

Perhaps one architecture diagram, one key results table, and one comparison of generated vs real images.

\section*{Conclusion}

The project implemented and evaluated a diffusion-based approach for human face generation using the CelebA-HQ dataset. The model was trained under the DDPM framework, where image generation is achieved by learning to iteratively denoise random Gaussian noise into realistic face images. Across the experiments, performance was evaluated using FID and supported by qualitative inspection of generated samples.

The results show that training duration has a major impact on sample quality. The baseline model improved substantially as training increased from 100 to 200 epochs, with FID decreasing from approximately 181 to 118. However, extending training to 250 epochs without additional stabilisation caused performance to degrade, increasing FID to around 159. This suggests that longer training alone is not sufficient, as the model can become unstable without techniques that smooth or regulate optimisation.

Exponential Moving Average was found to be an effective stabilisation method. EMA consistently improved sample quality compared to the raw model and allowed longer training runs to remain stable. With EMA, FID improved to approximately 95 at 200 epochs, 76 at 250 epochs, and 62 at 300 epochs. The strongest result under the restricted 2700-image training setup was achieved when EMA was combined with cosine learning rate scheduling. This configuration achieved an FID of 41.82 at 600 epochs, with visually clearer and more coherent generated faces.

A further scale-up experiment showed the practical importance of dataset size, batch size, and model capacity. Training on 81K images using 128×128 resolution, 128 base channels, batch size 16, and 100 epochs achieved an FID of 41.31. This was the best overall result and slightly outperformed the 600-epoch restricted-dataset model. Since dataset size, batch size, and model capacity were changed together, this result should be interpreted as evidence that scaling improves practical generation quality, rather than as an isolated comparison of any single factor.

Other experimental factors also influenced performance. Lower resolution images produced better FID scores in early experiments, but the generated samples lacked fine detail and realism. Therefore, 128×128 resolution was retained as a better trade-off between quantitative performance and visual quality. Changes in model capacity showed that both smaller and larger models could improve early results, but the 64-channel configuration provided the most reliable foundation for extended training under the restricted dataset setting. A lower learning rate led to slower convergence and inferior results compared to the baseline learning rate of 0.0002.

Despite these improvements, several limitations remain. Most experiments were conducted using only 2700 training images and a batch size of 1, which limited diversity, slowed training, and made optimisation noisier. Although the final models captured general facial structure, symmetry, and appearance, the generated images still do not fully match the fidelity of real samples. In addition, improvements beyond 450 epochs showed diminishing returns and some fluctuation, indicating that later-stage training remains difficult to stabilise.
 
Future work could improve the model by training consistently on larger datasets, using larger batch sizes, and exploring more advanced diffusion architectures. Additional improvements could include learned variance, classifier-free guidance, or higher-resolution training. More efficient sampling methods could also be investigated to reduce generation time. Overall, the experiments demonstrate that diffusion models are effective for human face generation, and that both stabilisation techniques and training scale are essential for achieving strong performance.

\section*{References}

\cite{b1}. 
\cite{b2}. 

Unless there are six authors or more give all authors' names; do not use 
``et al.''. Papers that have not been published, even if they have been 
submitted for publication, should be cited as ``unpublished'' \cite{b4}. Papers 
that have been accepted for publication should be cited as ``in press'' \cite{b5}. 
Capitalize only the first word in a paper title, except for proper nouns and 
element symbols.

\begin{thebibliography}{00}


\bibitem{b1} P. Dhariwal and A. Nichol, ``Diffusion models beat GANs on image synthesis,'' in Advances in Neural Information Processing Systems, vol. 34, pp. 8780–8794, 2021.
\bibitem{b2} Y. Song, J. Sohl-Dickstein, D. P. Kingma, A. Kumar, S. Ermon, and B. Poole, "Score-based generative modeling through stochastic differential equations," in Proceedings of the International Conference on Learning Representations (ICLR), 2021.
\bibitem{b3} Ho, J., Jain, A., \& Abbeel, P. (2020). Denoising Diffusion Probabilistic Models. NeurIPS 2020.
\bibitem{b4} J. Sohl-Dickstein, W. Eric, M. Niru, and G. Surya, "Deep Unsupervised Learning using Nonequilibrium Thermodynamics," in Proc. Int. Conf. Learn. Represent. (ICLR), 2015.
\bibitem{b5} M. Heusel, H. Ramsauer, T. Unterthiner, B. Nessler, and S. Hochreiter, "GANs trained by a two time-scale update rule converge to a local Nash equilibrium," in Advances in Neural Information Processing Systems, vol. 30, pp. 6626–6637, 2017.
\bibitem{b6} I. J. Goodfellow, J. Pouget-Abadie, M. Mirza, B. Xu, D. Warde-Farley, S. Ozair, A. Courville, and Y. Bengio, "Generative adversarial nets," in Advances in Neural Information Processing Systems, vol. 27, pp. 2672–2680, 2014.
\bibitem{b7} T. Karras, T. Aila, S. Laine, and J. Lehtinen, "Progressive Growing of GANs for Improved Quality, Stability, and Variation," in Proc. Int. Conf. Learn. Represent. (ICLR), 2018.
[8] A. Q. Nichol and P. Dhariwal, "Improved Denoising Diffusion Probabilistic Models," in Proc. Int. Conf. Mach. Learn. (ICML), pp. 8162–8171, 2021.
\bibitem{b8} A. Nichol and P. Dhariwal, "Improved Denoising Diffusion Probabilistic Models," in Proceedings of the International Conference on Machine Learning (ICML), pp. 8162–8171, 2021.
\bibitem{b9} L. Yang, Z. Zhang, Y. Song, S. Hong, R. Xu, Y. Zhao, W. Zhang, B. Cui, and M.-H. Yang, "Diffusion models: A comprehensive survey of methods and applications," ACM Computing Surveys, vol. 56, no. 4, pp. 1–39, 2024.

\end{thebibliography}
\vspace{12pt}
\color{red}
IEEE conference templates contain guidance text for composing and formatting conference papers. Please ensure that all template text is removed from your conference paper prior to submission to the conference. Failure to remove the template text from your paper may result in your paper not being published.

\end{document}
